\begin{apartats}

\apartat[1,25]{0,75}
Calcula la inversa de $A=\begin{pmatrix}5&3\\3&2\end{pmatrix}$ i comprova-la.

\begin{solucio}
$|A|=10-9=1\ne0$: és invertible. $A^{-1}=\dfrac1{|A|}(\operatorname{Adj}A)^t=\begin{pmatrix}2&-3\\-3&5\end{pmatrix}$.\\
Comprovació: $\begin{pmatrix}5&3\\3&2\end{pmatrix}\begin{pmatrix}2&-3\\-3&5\end{pmatrix}
=\begin{pmatrix}10-9&-15+15\\6-6&-9+10\end{pmatrix}=I$.
\end{solucio}

\begin{tria}{inversa-tasca}
\itemtria{ordre-3}{1}{1,25}
Calcula la inversa de $B=\begin{pmatrix}1&1&0\\0&1&1\\1&2&2\end{pmatrix}$.

\begin{solucio}
$|B|=2+1+0-0-2-0=1\ne0$: és invertible. Els adjunts de cada element són
\[
\operatorname{Adj}B=\begin{pmatrix}0&1&-1\\-2&2&-1\\1&-1&1\end{pmatrix},\qquad
B^{-1}=\frac1{|B|}(\operatorname{Adj}B)^t=\begin{pmatrix}0&-2&1\\1&2&-1\\-1&-1&1\end{pmatrix}.
\]
\end{solucio}

\itemtria{relaciona}{1}{1,25}
Relaciona cada matriu amb la seva inversa, sense calcular cap inversa:
\[
\text{a) }\begin{pmatrix}1&1\\0&2\end{pmatrix}\quad
\text{b) }\begin{pmatrix}2&1\\1&1\end{pmatrix}\quad
\text{c) }\begin{pmatrix}0&1\\1&2\end{pmatrix}\qquad
\text{I) }\begin{pmatrix}-2&1\\1&0\end{pmatrix}\quad
\text{II) }\frac12\begin{pmatrix}2&-1\\0&1\end{pmatrix}\quad
\text{III) }\begin{pmatrix}1&-1\\-1&2\end{pmatrix}
\]

\begin{solucio}
N'hi ha prou de comprovar quin producte dona la identitat. Per exemple,
$\begin{pmatrix}1&1\\0&2\end{pmatrix}\cdot\frac12\begin{pmatrix}2&-1\\0&1\end{pmatrix}
=\frac12\begin{pmatrix}2&0\\0&2\end{pmatrix}=I$.\\
Així: a) $\to$ II), b) $\to$ III) i c) $\to$ I). (També es poden reconèixer pel determinant: el de la
inversa és l'invers del de la matriu; a) té determinant 2, i II), $\frac12$.)
\end{solucio}
\end{tria}

\begin{nomesllarg}
\apartat{0,75}
Per què la matriu $\begin{pmatrix}2&4\\3&6\end{pmatrix}$ no té inversa? Troba el valor de $x$ per al qual
$\begin{pmatrix}2&4\\3&x\end{pmatrix}$ no és invertible.

\begin{solucio}
$\begin{vmatrix}2&4\\3&6\end{vmatrix}=12-12=0$: una matriu quadrada és invertible si i només si el seu
determinant no és nul.\\
$\begin{vmatrix}2&4\\3&x\end{vmatrix}=2x-12=0$ si $x=6$: només per a $x=6$ no és invertible.
\end{solucio}
\end{nomesllarg}

\end{apartats}
